\chapter{FIRMAS DE SEGUNDO ORDEN DE KOBIALKO Y DISCRIMINACIÓN DE MODELOS}

\section{Expansión de Taylor de Kobialko para la sombra \texorpdfstring{$R^2(\delta)$}{R^2(delta)}}
Siguiendo a Kobialko \& Gal'tsov (2024) \citep{perturbaciones_en_bh}, expandimos el cuadrado del parámetro de impacto crítico $b_c^2$ hasta segundo orden en el parámetro de carga débil $\delta$:
\begin{equation}
    b_c^2(\delta) = 27 r_g^2 \left( 1 - B \delta - D \delta^2 + \mathcal{O}(\delta^3) \right)
    \label{eq:kobialko_expansion_SI}
\end{equation}
donde $r_g = \frac{GM}{c^2}$. Para la métrica de Kerr-Newman, el coeficiente lineal es $B = 1/3$, idéntico al de otros modelos cargados \citep{perturbaciones_en_bh}.

\section{Aislamiento del coeficiente de segundo orden \texorpdfstring{$D = 1/27$}{D = 1/27}}
Deducimos paso a paso el coeficiente de segundo orden para Kerr-Newman cargado débilmente \citep{perturbaciones_en_bh}:
\begin{equation}
    D_{KN} = \frac{1}{27}
\end{equation}
Comparando analíticamente con modelos alternativos \citep{perturbaciones_en_bh}:
\begin{equation}
    D_{Bardeen} = \frac{2}{27}, \quad D_{Hayward} = \frac{8}{243}
\end{equation}
El coeficiente $D = 1/27$ es la firma única e inconfundible de la carga eléctrica de electrovacío en segundo orden.

\section{El observador geodésico y aberración relativista}
El marco local del observador geodésico introduce correcciones por aberración relativista en las coordenadas celestes $(\alpha, \beta)$ \citep{shadows_quintessence_bh, Achour_2025_J._Cosmol._Astropart._Phys._2025_012}, transformando el perfil angular medido.
